\documentclass[10pt,a4paper]{amsart}
\usepackage[margin=19mm]{geometry}
\usepackage{amsmath,amssymb,mathtools}
\usepackage[scr=rsfs,cal=euler]{mathalfa}
\usepackage{xfrac}
\usepackage[unicode,hidelinks,bookmarks=false]{hyperref}
\newcommand{\ie}{i.e.\ }
\newcommand{\cf}{cf.\ }
\newcommand{\resp}{respectively\ }
\newcommand{\Subsection}{\subsection}
\providecommand{\nobreakdash}{\nobreak\hbox{-}}
\setlength{\emergencystretch}{2em}
\multlinegap0pt
\usepackage{amssymb}
\usepackage{mathtools}
\usepackage[shortlabels]{enumitem}
\newtheorem{theorem}{Theorem}
\usepackage{cleveref}
\crefname{theorem}{Theorem}{Theorems}
\newcommand{\BCTWO}{\mathsf{BCT}_{\WO}}
\newcommand{\R}{\mathbb R}
\newcommand{\TP}{\mathsf{TP}_{\WO}}
\newcommand{\WO}{\mathsf{WO}}
\title{Ideals, Well-Orderable Families, and Baire Category in Truss's Feferman-Type Model}
\author{Jason Zesheng Chen}
\date{Source: arXiv:2608.21649v1 (CC BY 4.0)}
\begin{document}
\maketitle

\noindent\emph{Source and licence.} Ideals, Well-Orderable Families, and Baire Category in Truss's Feferman-Type Model. Version arXiv:2608.21649v1 (CC BY 4.0), distributed by arXiv under the Creative Commons Attribution 4.0 International licence, http://creativecommons.org/licenses/by/4.0/, as recorded in the arXiv licence metadata for this exact version and shown on the abstract page https://arxiv.org/abs/2608.21649v1. Full licence notice: \texttt{licenses/arxiv-2608.21649-CC-BY-4.0.txt}.\par\smallskip

\noindent\textit{Editorial scope.} Counted unit: the article's theorem thm:model-principles (source0.tex lines 523--528). The article's complete original proof of it is delivered with it (source0.tex lines 530--550, one \texttt{proof} environment of 21 lines). No mathematics is written, repaired or re-derived: the statement and the proof are the article's own lines, in the article's own notation, and no retained line is substituted or re-flowed.  Retained uncounted as the source-local context the proof needs: lines 454--467.  Nothing was omitted from any retained span, so the package carries no omission record.  No retained line contains a citation, so the wrapper carries no bibliography and no bibliography entry.  No retained line contains a figure environment, an image-inclusion command or drawing code, so no image is delivered and the package declares no asset dependency.  Editorial formatting only, in addition: an AMS \texttt{amsart} preamble that restates the article's own declarations and macros used by the retained lines, namely the environments \texttt{theorem} on separate counters, together with the standard packages those lines need.  The retained environments print in this document as Theorem 1 in order of appearance, whereas the article prints the same items with the numbering of its own full document.  The complete native source of the article as distributed in the arXiv e-print for this exact version is bundled unchanged as \texttt{sources/208264/source0.tex}.
\begin{theorem}\label{thm:truss}
The model $M$ has the following properties.
\begin{enumerate}
\item $M\models\mathsf{ZF}+\mathsf{AC}_{\WO}$.
\item If $A\subseteq\R$ belongs to $M$, then
\[
A\text{ is well-orderable in }M
\quad\Longleftrightarrow\quad
\exists\alpha<\omega_1\quad A\subseteq\R^{L[G_{<\alpha}]}.
\]
\item Every non-well-orderable set of reals contains a perfect subset.
\item $M\models V=L(\R)$, and hence $M\models\mathsf{SVC}(\R)$.
\end{enumerate}
\end{theorem}

\begin{theorem}\label{thm:model-principles}
The model $M$ satisfies
\[
\mathsf{ZF}+\mathsf{DC}+\mathsf{AC}_{\WO}+\TP+\BCTWO.
\]
\end{theorem}

\begin{proof}
Only $\BCTWO$ remains to be proved.
Let $X$ be a perfect Polish space in $M$, and let $\mathcal D$ be a well-orderable family of dense open subsets of $X$.
Fix a complete compatible metric and a countable dense sequence for $X$, and let $\langle U_n:n<\omega\rangle$ be the resulting countable basis.
For $D\in\mathcal D$, code $D$ by
\[
d_D=\{n<\omega:U_n\subseteq D\}.
\]
The set $\{d_D:D\in\mathcal D\}$ is a well-orderable set of reals, so by \Cref{thm:truss} there is $\alpha<\omega_1$ such that every $d_D$ belongs to $L[G_{<\alpha}]$.
Increasing $\alpha$ if necessary, we may also assume that the chosen presentation of $X$ belongs to $L[G_{<\alpha}]$.

Let $U$ be a nonempty basic open subset of $X$.
Over $L[G_{<\alpha}]$, the forcing of nonempty basic open subsets of $U$ ordered by reverse inclusion has a countable atomless dense suborder, and hence is forcing-equivalent to Cohen forcing.
The fresh coordinate $c_\alpha$ therefore yields a point $x\in U$ which is Cohen-generic for the topology of $X$ over $L[G_{<\alpha}]$.
Since each $d_D$ belongs to $L[G_{<\alpha}]$ and codes a dense open set there, $x\in D$ for every $D\in\mathcal D$.
Thus
\[
U\cap\bigcap\mathcal D\neq\varnothing.
\]
Since $U$ was arbitrary, $\bigcap\mathcal D$ is dense in $X$.
\end{proof}

\end{document}
